% NeurIPS Paper Checklist (mandatory)
% Do not remove this section.

\section*{NeurIPS Paper Checklist}

\begin{enumerate}

\item {\bf Claims}
    \item[] Question: Do the main claims made in the abstract and introduction accurately reflect the paper's contributions and scope?
    \item[] Answer: \answerYes{}
    \item[] Justification: The abstract and introduction state four specific contributions (Section~\ref{sec:intro}), each supported by empirical evidence in Sections~\ref{sec:results:clinicalbench}--\ref{sec:results:slicebench}. All claims include confidence intervals and are scoped with ``to our knowledge'' qualifiers where appropriate.

\item {\bf Limitations}
    \item[] Question: Does the paper discuss the limitations of the work performed by the authors?
    \item[] Answer: \answerYes{}
    \item[] Justification: Section~\ref{sec:discussion} explicitly discusses limitations including small sample sizes (6 patients for SliceBench, 2 per tier), single-site data (MIMIC-IV), the B2$\rightarrow$B3 confidence interval crossing zero, category-level trade-offs (C4 vs C1), and the ``change'' category capability gap.

\item {\bf Theory assumptions and proofs}
    \item[] Question: For each theoretical result, does the paper provide the full set of assumptions and a complete (correct) proof?
    \item[] Answer: \answerNA{}
    \item[] Justification: This is a systems paper with empirical evaluation; no theoretical proofs are claimed.

\item {\bf Experimental result reproducibility}
    \item[] Question: Does the paper fully disclose all the information needed to reproduce the main experimental results of the paper to the extent that it affects the main claims and/or conclusions of the paper?
    \item[] Answer: \answerYes{}
    \item[] Justification: Appendix~\ref{app:reproducibility} reports exact model versions (MedGemma 27B, Claude Sonnet 4.5, Claude Opus 4.6), temperature settings (0), bootstrap parameters ($n=2{,}000$, seed 42, BCa method), and data source (MIMIC-IV v1.5.0 under PhysioNet DUA). Section~\ref{sec:benchmarks} describes question generation and evaluation protocols.

\item {\bf Open access to data and code}
    \item[] Question: Does the paper provide open access to the data and code, with sufficient instructions to faithfully reproduce the main experimental results?
    \item[] Answer: \answerNo{}
    \item[] Justification: MIMIC-IV data is available under PhysioNet credentialed access (not open). Code will be released upon acceptance. Benchmark question templates and evaluation scripts will be included in the release.

\item {\bf Experimental setting/details}
    \item[] Question: Does the paper specify all the training and test details necessary to understand the results?
    \item[] Answer: \answerYes{}
    \item[] Justification: Section~\ref{sec:benchmarks} specifies benchmark design, question counts, condition definitions, patient selection criteria, and evaluation methodology. Appendix~\ref{app:reproducibility} provides model versions and computational requirements.

\item {\bf Experiment statistical significance}
    \item[] Question: Does the paper report error bars suitably and correctly?
    \item[] Answer: \answerYes{}
    \item[] Justification: Table~\ref{tab:slicebench_overall} reports BCa bootstrap 95\% CIs for all conditions. Per-tier results (Figure~\ref{fig:slicebench_figures}) are clearly labeled as point estimates without per-tier CIs due to small tier-level sample sizes ($n=2$ patients per tier). The overall B2$\rightarrow$B3 CI crossing zero is explicitly noted.

\item {\bf Experiments compute resources}
    \item[] Question: For each experiment, does the paper provide sufficient information on the computer resources needed to reproduce the experiments?
    \item[] Answer: \answerYes{}
    \item[] Justification: Appendix~\ref{app:reproducibility} reports computational requirements: ClinicalBench $\approx$3.6 hours (single GPU for MedGemma inference), SliceBench $\approx$2 hours (Anthropic API).

\item {\bf Code of ethics}
    \item[] Question: Does the research conducted in the paper conform, in every respect, with the NeurIPS Code of Ethics?
    \item[] Answer: \answerYes{}
    \item[] Justification: The study uses de-identified clinical data under an approved data use agreement (Appendix~\ref{app:ethics}). No patient re-identification was attempted. The system is designed for clinical decision \emph{support}, not autonomous decision-making.

\item {\bf Broader impacts}
    \item[] Question: Does the paper discuss both potential positive societal impacts and negative societal impacts of the work?
    \item[] Answer: \answerYes{}
    \item[] Justification: Section~\ref{sec:discussion} discusses the potential for improved clinical decision support (positive) and risks of over-reliance on automated systems, context dilution, and the need for physician oversight (negative). Appendix~\ref{app:ethics} addresses data ethics.

\item {\bf Safeguards}
    \item[] Question: Does the paper describe safeguards that have been put in place for responsible release of data or models?
    \item[] Answer: \answerYes{}
    \item[] Justification: The system operates on de-identified data under PhysioNet DUA. It is designed as a decision support tool requiring physician oversight. Code release will include usage guidelines emphasizing the system is not validated for autonomous clinical use.

\item {\bf Licenses for existing assets}
    \item[] Question: Are the creators or original owners of assets used in the paper properly credited and are the license and terms of use explicitly mentioned and properly respected?
    \item[] Answer: \answerYes{}
    \item[] Justification: MIMIC-IV is cited~\cite{johnson2023mimiciv} and used under PhysioNet Credentialed Health Data Use Agreement. All benchmark frameworks and models are cited. OMOP CDM is an open standard.

\item {\bf New assets}
    \item[] Question: Are new assets introduced in the paper well documented and is the documentation provided alongside the assets?
    \item[] Answer: \answerYes{}
    \item[] Justification: ClinicalBench and SliceBench are new benchmarks fully documented in Section~\ref{sec:benchmarks}, with question generation procedures, category definitions (Appendix~\ref{app:assertions}), and evaluation protocols. Templates and scoring scripts will accompany the code release.

\item {\bf Crowdsourcing and research with human subjects}
    \item[] Question: For crowdsourcing experiments and research with human subjects, does the paper include the full text of instructions given to participants and screenshots?
    \item[] Answer: \answerNA{}
    \item[] Justification: No crowdsourcing was used. Physician adjudication (Appendix~\ref{app:human}) involves expert review by co-authors, not human subjects research.

\item {\bf IRB approvals or equivalent}
    \item[] Question: Does the paper describe potential risks and does it include the IRB approval number if applicable?
    \item[] Answer: \answerNA{}
    \item[] Justification: MIMIC-IV is a de-identified public dataset; use under PhysioNet DUA does not require separate IRB approval per PhysioNet policy.

\end{enumerate}
